% !TeX root = ../papers/nc-p3.tex
\section{三次曲面上の階数二曲線と核の鉛筆}
\label{sec:cofactor-cubic-surface}
\begin{theorem}\label{thm:cofactor-no-integral-cubic}
被約四辺有限窓の一般可逆な全階数反対根で、余因子スキームの曲線部分全体が一般被約かつ
その台 $C$ が整なら、$C$ を含む三次曲面は存在しない。
\end{theorem}
\begin{proof}[一般平面三次上の二つの線形系]
存在すれば定理\ref{thm:cofactor-no-quadrics}よりその曲面 $F$ は整である。同定理から
$h^0(\omega_C(-3))\ge2$。$J=\mathcal I_{C/F}$ と置くと
\[
 H^0\bigl(F,\mathcal Hom_F(J,\mathcal O_F(-4))\bigr)
                      \simeq H^0(C,\omega_C(-3)).
\]
実際 $0\to J\to\mathcal O_F\to\mathcal O_C\to0$ を $\omega_F(-3)=\mathcal O_F(-4)$ で
双対化し、余次元一の双対化層と $H^0(\mathcal O_F(-4))=H^1(\mathcal O_F(-4))=0$ を使う。
$F$ の正規性、$C$ のCartier性は不要である。
独立な二切断の一般点での比は非定数であり、一般平面切断 $E$ に制限しても非定数である。
実際一般二点を通る平面の入射を使えば、その比が一般平面切断で定数なら一般二点の値が等しくなる。
$E$ は整な算術種数一の平面三次で、$C$ の $d$ 点は $F$ の滑らかな点で横断的に切れる。
$C$ が曲面の特異軌跡に含まれると三次の非零偏導二次式が $C$ を含むためである。
従って二切断は $E$ 上で次数 $d-12$ の同じ線束の独立な切断になる。
整な平面三次では次数一以下の線束の切断は高々一つである（双対化層は自明でRiemann--Rochを使う）。
よって $d\ge14$ である。

$E$ 上で行列式の次数は二十四、各交点で少なくとも二重なので行列式は $F$ で消える。
階数二の曲面はなく、$F$ の一般階数は三である。一般の $E$ では一般核点の射が $E$ 全体に延長する。
実際 $E$ の特異点でも $C$ と孤立した階数二以下の点を避けるので階数三であり、
$C$ 上の $E$ の点は滑らかである。この射の線束を $L_E$ とすれば、
余因子はそのVeronese座標に共通切断を掛けたもので、共通切断は
$\mathcal O_E(6)\otimes L_E^{-2}$ の切断として $d$ 点で消える。
従って $2\deg L_E\le18-d\le4$。$F$ 全体で核一定なら固定列が三次式で割れて矛盾するので、
一般 $E$ 上でも射は非定数。再び算術種数一により $\deg L_E=2$, $h^0(L_E)=2$ で、
$d=14$、各 $E$ の核像は直線である。この直線は $E$ に依らない。
実際一般の二平面の共通直線上の三点のうち少なくとも二点では核が異なる。
一般二点を $F$ 上で選びその割線を取る写像は一般直線へ支配的かつ生成的有限
（一般直線は相異なる三点を切る）であり、二点の核は一般に異なるからである。
二平面の核直線はその両方を通るので一致する。
\end{proof}

\begin{proof}[核が一直線を動く三次因子]
このような既約三次因子 $F$（方程式も $F$ と書く）があれば、
$N$ の階数二以下の軌跡は必ず一本の直線を含むことを示す。
一般核を含む二空間の基底を $v,w$ とすると、
\[
             Nv\wedge Nw=F\,D(x)
\]
で $D$ は線形な六座標、一般に非零である。$D$ の射影像の閉包はGrassmann多様体内の線形空間だから、
全二空間がある固定三空間に含まれるか、全二空間がある固定直線 $kz$ を含む。
これはPlücker直線上の二平面が交わること、および二本の交わる射影直線と両方交わる直線は
その交点を通るかその張る平面に入ることからも従う。
前者なら $Nv$ の四成分が独立でなくなるので、後者を取る。
$z$ で割った二列は比例する。多項式環で共通の原始三成分列をくくれば
\[
             \overline{Nv}=l_1 b,\qquad \overline{Nw}=l_2 b
\]
と書ける。定数係数の二列の関係は非零の一列の像を $kz$ に入れてしまい、
原始列が定数でも各列の三射影成分の独立性が失われる。
従って $l_1,l_2$ は独立な一次式、$b$ の三成分も独立な一次式である。
$N(l_2v-l_1w)=h z$ の非零三次式 $h$ は $F$ の一般点で消える：
そこでは二列が比例し、その射影が同時には零でないからである。従って $h$ は $F$ の定数倍である。
三次評価が単射なので $F(p)\ne0$ となる印があり、そこで
$N(p)$ は $z$ の平方に比例し、$z$ は $\langle v,w\rangle$ を消さない。
基底を選び直し、最初の二つを $v,w$ とし、双対の第一ベクトルを $z$ にできる。
射影二列の分解はこの基底でも成立する。対称性により表示は
\[
 N=\begin{pmatrix}
 a & l_1b_1 & l_1b_2 & l_1b_3\\
 l_1b_1 & l_2b_1 & l_2b_2 & l_2b_3\\
 l_1b_2 & l_2b_2 & c & e\\
 l_1b_3 & l_2b_3 & e & g
 \end{pmatrix},\qquad
             F\sim a l_2-l_1^2 b_1 .
\]
特に $b_1$ は $l_2$ に比例しない。三式
\[
         b_1c-l_2b_2^2,\quad b_1e-l_2b_2b_3,\quad b_1g-l_2b_3^2
\]
は線形独立である。関係を $b_1$ を法として読むと、$b_2,b_3$ はその面で独立であり、
$l_2$ は零でないからである。印でこれらが消えない可能性があるのは
$l_1=l_2=a=0$ の点に限られる。実際 $l_2\ne0$ では二行小行列式の消滅を $l_2$ で割ればよい。
$l_2=0$, $b_1\ne0$ なら階数一の第二対角成分が零なので第二列も零、従って $l_1=0$。
さらに $a\ne0$ なら第一列が $az$ で出力点は $z$、$c=e=g=0$ である。
三次評価の単射性から三式を支える相異なる印が三点以上必要で、
$a$ は直線 $l_1=l_2=0$ において恒等零となる。この直線上では最初の二列とも消え、
主張した階数二以下の直線を得る。整な $C$ がこの直線を含むことは
定理\ref{thm:cofactor-no-quadrics}に反する。
\end{proof}
